%Texlive-full Version 3.141592-1.40.3 (Web2C 7.5.6)
%Kile Version 2.0.83

%File associated : Mas_Logo.ps , uncertainty.bib

\documentclass[a4paper,10pt]{article}
\usepackage[utf8x]{inputenc}

\usepackage{lmodern}
\usepackage[a4paper]{geometry}
%\usepackage[frenchb]{babel}
\usepackage{graphicx}
\usepackage{hyperref}

\usepackage{pstricks}
\usepackage{pst-node}
%\usepackage{wrapfig}
\usepackage{amsmath}
\usepackage{amssymb}

\usepackage{listings}
\lstset{language=C++,basicstyle=\scriptsize \color{green},identifierstyle=\color{orange},keywordstyle=[1]\color{blue},columns=fullflexible}

\begin{document}
%%%%%%%%%%%%%%%%%%   LOGO  %%%%%%%%%%%%%%%%%%%%%%%%%
\begin{center}
\rput(6,1.5){\href{http://www.mas.ecp.fr/}{\includegraphics[scale=1.0]{Mas_Logo}}}
\rput(-4,1.5){\href{http://www.mas.ecp.fr/}{
		\begin{tabular}{l}
		\resizebox{3cm}{0.6cm}{PyGMalion} \\ 
		\resizebox{6cm}{0.3cm}{Digiplant Tech.Doc}
		\end{tabular}
		}
	    }
\end{center}
%%%%%%%%%%%%%%%%%%   LOGO  %%%%%%%%%%%%%%%%%%%%%%%%%
%%%%%%%%%%%%%%%%%% DOCUMENT TITLE %%%%%%%%%%%%%%%%%%%%%%%%%
\vspace{1.5cm}
\begin{center}\resizebox{5cm}{0.6cm}{Uncertainty Analysis}\end{center}
%%%%%%%%%%%%%%%%%% DOCUMENT TITLE %%%%%%%%%%%%%%%%%%%%%%%%%
\section{Functionalities }
\subsection{Uncertainty Analisis } 
%%%%%%%%%%%%%%%%%%%%%%%%%%%%%%%%%%%%%%%%%%%%%%%%%%%%%%%%%%%%%%%%%%%%%%%%%%%%%%%%%%%%%%%%%
%=======================================================================================%
%%%%%%%%%%%%%%%%%%%%%%%%%%%%%%%%%%%%%%%%%%%%%%%%%%%%%%%%%%%%%%%%%%%%%%%%%%%%%%%%%%%%%%%%%
\subsubsection{Unscented Tools: } 
\paragraph{Abstract : } The Unscented Tool work by associating to a function
linear or not $f$.
%%%%%%%%%%%%%%%%%%%%%%%%%%%%%%%%%%%%%%%%%%%%%%%%%%%%%%%%%%%%%%%%%%%%%%%%%%%%%%%%%%%%%%%%%
%============================           FIGURE        ==================================%
\[
f: 
\left\{ 
\begin{array}{lll}
     \textbf{R}^{N} &\rightarrow&  \textbf{R}^{M}   \\
     \textbf{X}     &\mapsto&      \textbf{Y} 
\end{array}\right.
\]
\begin{figure}[h]
\begin{center}
\begin{pspicture}(-5,-2)(5,2)
%\psline(-5,-2)(5,2)
\rput(-4, 1){\Rnode{mX}{$\{\bar{\textbf{X}},\textbf{P}_{xx}\}$}}
\rput(-4,-1){\Rnode{Xi}{$\textbf{\Large{X}}_i$}}
\rput( 0, 0){\Rnode{F}{$\textbf{\Large{f}}$}}
\rput( 4,-1){\Rnode{Yi}{$\textbf{\Large{Y}}_i$}}
\rput( 4, 1){\Rnode{mY}{$\{\bar{\textbf{Y}},\textbf{P}_{yy}\}$}}
\psellipticarc{<-}(0,0)(0.5,1){60}{320}\rput(0.5,1){\Large{i}}
%%%%%%%%%%%%%%%%%%%%%%%%%%%%%%%%%%%%%%%%%%%%%%%%%%%%
\ncline{->}{mX}{Xi}\aput{:U}{unscented}
\ncline{->}{Yi}{mY}\aput{:U}{unscented}
\ncline[linestyle=dotted]{->}{Xi}{F}
\ncline[linestyle=dotted]{->}{F}{Yi}
\end{pspicture}
\end{center}
\end{figure}
%============================           FIGURE        ==================================%
%%%%%%%%%%%%%%%%%%%%%%%%%%%%%%%%%%%%%%%%%%%%%%%%%%%%%%%%%%%%%%%%%%%%%%%%%%%%%%%%%%%%%%%%%
In general way, Unscented Tool may be used for many case : for example Unscented Kalman Filter or the Unscented Transform method in uncertainty analysis. For the reusability, the general senarios workflow of UnscentedTool composed steps follows : 
\begin{enumerate}
 \item given mean and covariance of a random vector$\{\bar{\textbf{X}},\textbf{P}_{xx}\}$(dimension \textbf{N})
 \item generate \textbf{2.N+1} point $\textbf{X}_{i,0\leq i< (2N+1)}$
 \item \ldots doing the stuff for getting $\textbf{Y}_{i} = f(\textbf{X}_{i}), _{(i,0\leq i< (2N+1))}$\footnote{This step this done by other component. Unscented Tool doesn't do it}
 \item compute the mean and the covariance of transformed points $\{\bar{\textbf{Y}},\textbf{P}_{yy}\}$
\end{enumerate}
Several variants of step 2 and step 4 are publied by Julier et al. The first one basical is on 1996~\cite{Julier1996}, the advanced one introduced a scaling factor at 2000~\cite{Julier2002} allowing when there a large angular uncertainty (non-diagonal term of cov matrix?).   
\paragraph{Implementation : } An UnscentedSTool is simply a set of weights.
These weights allows computation of step 2 and steps 4 above. By Julier, all
weight values can be computed by giving tree factor $\kappa$ ,$\alpha$ ,
$\beta$. These weights are not much different, and not need to be computed for all in the mathematical point of view. But in the programming point of view, when step 2 and steps 4 need to be done repeatly, it may evoke a hight cost of time if weights are recompute every time. Therefor weights are hiddenly all precomputed and stored in the components Unscented Tool. When user call for step 2 or step 4, it apply these weights to points and compute :
\begin{itemize}
 \item  step 1 : \textbf{compSigmas}:
      \[
      \left\| 
      \begin{array}{l}
	    \text{for : } 0\leq j < \textbf{N} \textbf{ : } \sigma_j = \textbf{sqrt}(\textbf{P}_{xx}).\text{col}_{(j)}\\
	    \textbf{X}_0 = \bar{\textbf{X}} \\
	    \text{for : } 1\leq i < \textbf{2N+1} \textbf{ : } \textbf{X}_i = \bar{\textbf{X}}+ \textbf{w}^{x}_{i}*\sigma_i
      \end{array}\right.
      \]
 \item  step 2a : \textbf{compMean}:
      \[
      \left\| 
      \begin{array}{l}
	    \text{for : } 0\leq i < \textbf{2N+1} \textbf{ : } \bar{\textbf{Y}} += \textbf{w}^{y}_{i}*\textbf{Y}_i
      \end{array}\right.
      \]
 \item  step 2b : \textbf{compCov}:
      \[
      \left\| 
      \begin{array}{l}
	    \text{for : } 0\leq i < \textbf{2N+1} \textbf{ : } \textbf{P}^{lm}_{yy} += \textbf{w}^{yy}_{i}*\left[\textbf{Y}_i - \bar{\textbf{Y}} \right]^{l}  * \left[\textbf{Y}_i - \bar{\textbf{Y}} \right]^{m}
      \end{array}\right.
      \]
\end{itemize}
\subsubsection{Test : }
\paragraph{For linear case : }
When $f$ is an linear transform, there are an exact uncertainty transform through $f$. That mean after the unscented linear transform, we must have :
\[
\left\{ 
\begin{array}{ccl}
     \textbf{P}_{yy}  &=& \left[f\right]\textbf{P}_{xx}\left[f\right]^t  \\
     \bar{\textbf{Y}} &=& \left[f\right]\bar{\textbf{X}} 
\end{array}\right.
\]
\paragraph{Tolerance error : }
The error can provided from Cholesky factorization, making that result computed are not absolutely equal to the mathematical result~\cite{ErrorTest1}. Test tools is used of boost library~\cite{BoostCheck}. For instance BOOST\_CHECK\_SMALL (diff, tolerance) seems good. 
\paragraph{Keyword: } Covariance Matrix, Gaussian Distribution, Unscented Kalman Filter. 

						      %%%%%%%%%%%%%%%%%%%%%%%%%%  Writer %%%%%%%%%%%%%%%%%%%%%%%%
						      \begin{flushright}
						      Document written by \\
						      \href{mailto:chi-thanh.nguyen@inria.fr}{{\textbf {Chi Thanh NGUYEN}}} \\
						      MAS Central Paris  
						      \end{flushright}
						      %%%%%%%%%%%%%%%%%%%%%%%%%%  Writer %%%%%%%%%%%%%%%%%%%%%%%%
\bibliographystyle{siam}
\bibliography{uncertainty}
\end{document}